\chapter{绪论}\label{chap:intro}

\section{研究背景及意义}

以大语言模型为代表的生成式人工智能正加速渗透到内容创作、代码生成、智能问答等生产生活的各个环节，其生成内容的安全性直接关系到公众利益与社会秩序。与此同时，模型权重的开放发布在推动技术普惠的同时，也使有害知识的治理面临更大挑战。我国《生成式人工智能服务管理暂行办法》明确要求生成式人工智能服务提供者承担内容安全主体责任，对违法内容履行处置义务；如何高效治理大模型中的有害知识、保障生成内容安全，已成为学术界与产业界共同关注的核心议题。

从技术路线上看，现有的大模型内容安全治理手段可分为三类。其一为对齐训练，通过基于人类反馈的强化学习\cite{ouyang2022instruct}或直接偏好优化\cite{rafailov2023dpo}等方式在训练阶段注入安全偏好，此类方法依赖较大规模的标注数据与计算资源，且属于对模型的全局干预。其二为推理时过滤，通过外挂的护栏模型或安全分类器在推理阶段拦截有害输出，此类方法不改变模型本身，但引入推理开销，且难以应对经过伪装的越狱攻击。其三为知识编辑，通过对模型参数的定点修改，从模型内部精准擦除或修正有害知识\cite{zhang2024comprehensive}——相比前两者，知识编辑无需重训练、无需外挂组件、编辑效果直接固化于模型参数，为内容安全治理提供了一种轻量且精准的全新手段。三类技术路线的对比见图~\ref{fig:chap1-route}。

% TODO 图1-1：绘制后删除本占位框，替换为 \includegraphics[width=0.9\textwidth]{figures/chap01-fig1-route.pdf}
\begin{figure}[htbp]
  \centering
  \fbox{\parbox[c][6.5cm][c]{0.82\textwidth}{\centering\large 图1-1 占位\\[0.6em] 大模型内容安全治理的三类技术路线对比\\[0.4em] \small （对齐训练 / 推理时过滤 / 知识编辑：作用阶段与优缺点对比）}}
  \caption{大模型内容安全治理的三类技术路线对比}
  \label{fig:chap1-route}
\end{figure}

然而，将知识编辑应用于内容安全场景时，两个关键缺陷逐渐显现。其一，编辑不持久：已有研究表明，编辑知识在后续微调中会显著衰减\cite{cheng2025finetune}，且少量有害样本的微调即可破坏模型的安全对齐\cite{qi2023finetuning}；当被擦除的有害知识随常规微调而回潮时，安全编辑的失效将直接构成安全事故。其二，编辑不分级：现有方法以单一强度实施编辑，对暴力威胁与冒犯性言论等风险等级迥异的内容无差别处置，强度过高诱发对合法查询的过度拒绝，强度不足则放任高危内容\cite{shi2025cku,lu2025toxedit}。这两个缺陷分别从可靠性与可控性两个维度制约着知识编辑在内容安全场景的实际应用，也构成了本文的研究出发点。

在理论层面，本文面向上述缺陷开展的知识编辑持久性机制与分级可控编辑研究，将填补参数编辑范式在内容安全应用中的方法空白，为构建一次编辑、持久生效、按需调控的安全编辑方法体系提供支撑。在应用层面，研究成果可为开源大模型的有害内容治理提供轻量化的技术工具链，配套开发的可控知识编辑防御系统能够支撑安全策略的灵活配置与编辑效果的实际验证，对推动生成式人工智能的安全可信部署具有现实意义。

\section{国内外研究现状}

\subsection{大模型知识编辑技术}

大语言模型在预训练中习得海量世界知识，但其中不可避免地包含过时、错误乃至有害的内容，而通过全量重训练更新知识的计算成本难以承受。知识编辑旨在不显著影响模型整体能力的前提下，高效修正、注入或删除模型中的特定知识\cite{zhang2024comprehensive,wang2024survey}，涵盖知识插入、知识修改与知识删除三类基本场景，理想的编辑应满足目标知识被改变而无关知识不受影响的局部性约束。

按照知识与模型的结合方式，现有编辑方法可划分为三类\cite{zhang2024comprehensive}。第一类借助外部知识，通过检索增强\cite{lewis2020rag}或上下文学习使模型感知待更新的知识，此类方法无需修改参数，但知识并未真正写入模型，且引入额外的推理开销与外部依赖。第二类将新知识整合进模型，以参数高效微调\cite{hu2022lora}与元学习为代表，如MEND\cite{mitchell2022mend}通过超网络预测参数修正量以实现快速编辑。第三类直接编辑模型的内在知识，以定位-编辑范式为代表：ROME\cite{meng2022rome}与MEMIT\cite{meng2023memit}通过因果追踪定位事实知识存储的关键参数并直接修改权重，将编辑效果永久固化于模型参数之中。

定位-编辑范式建立在事实知识主要存储于前馈层的观察之上：ROME以闭式解修改单个事实对应的关键参数，MEMIT将其扩展至批量编辑场景，AlphaEdit\cite{fang2024alphaedit}进一步将编辑扰动投影至保留知识的零空间，实现编辑对原有知识的零干扰，该工作获ICLR 2025杰出论文奖，代表了该范式的最新进展。在评估方面，编辑成功率、可移植性、局部性与流畅性构成四项通用指标\cite{zhang2024comprehensive}；KnowEdit等基准与EasyEdit\cite{wang2024easyedit}开源框架为方法的实现与比较提供了统一平台。参数级编辑无需附加模块与推理开销，是当前兼顾效率与持久性的主流路线，也是本文采用的技术基础。

然而，参数级编辑将编辑效果一次性写入参数，其可靠性高度依赖参数自身的稳定性。当编辑对象从普通事实扩展到有害知识、部署环境从静态评测转向持续微调时，参数稳定性问题将直接威胁安全应用的可靠性，相关研究将在1.2.3节详细展开。

\subsection{大模型生成内容安全}

大语言模型的生成内容面临多维度安全风险，包括暴力、仇恨、自残等有害内容生成，偏见与歧视，以及通过角色扮演、指令重构等手段诱导模型突破安全约束的越狱攻击。早期的安全治理以对齐训练为主，通过基于人类反馈的强化学习\cite{ouyang2022instruct}与直接偏好优化\cite{rafailov2023dpo}等方式使模型习得拒绝有害请求的行为；然而对齐训练依赖高质量的偏好数据与可观的训练资源，且其安全效果容易被少量有害样本的微调所破坏\cite{qi2023finetuning}。

为从模型内部消除有害知识，机器遗忘技术被引入安全场景。WMDP基准\cite{li2024wmdp}系统评测了模型对生物、网络等危险知识的掌握程度，并提出表征误导向遗忘方法（RMU），在不显著损害通用能力的前提下移除危险知识；NPO\cite{zhang2024npo}从偏好优化视角重构遗忘目标，缓解梯度上升导致的灾难性能力崩溃；CKU\cite{shi2025cku}通过对多层感知机层的神经元评分，在遗忘有害知识的同时保护有用知识。遗忘类方法将安全治理从行为层面推进到知识层面，但其训练式本质决定了需要完整的目标函数设计，且普遍采用一刀切的遗忘强度。

知识编辑则为内容安全提供了另一条治理路径，相关防御工作近两年快速兴起。Wang等人\cite{wang2024safeedit}构建了覆盖九类有害行为的SafeEdit基准，并提出仅用单条样本定位并修改毒性区域的DINM方法，验证了知识编辑在减毒效率与通用能力保持上的潜力；ToxEdit\cite{lu2025toxedit}在前向传播中动态检测毒性激活模式并自适应路由计算路径，在精准减毒的同时抑制过度编辑；DELMAN\cite{wang2025delman}将模型编辑用于越狱防御，仅更新与越狱行为相关的最小参数集，并可动态适应新型攻击模式。这一系列工作表明，知识编辑正在成为内容安全治理中兼具效率与精准性的新选项。

然而，编辑类安全防御整体仍处于起步阶段：一方面，现有工作的评估多止于编辑后的即效性，安全编辑在持续微调部署环境下的持久性与恶意微调下的鲁棒性尚未进入研究视野；另一方面，防御对象多集中于毒性或越狱等单一风险类别，编辑强度普遍一刀切，缺乏按内容风险等级分级控制的手段。这两方面的局限将结合编辑持久性与安全风险的研究进展，在1.2.3节进一步展开。

\subsection{知识编辑的持久性与安全风险}

知识编辑以局部参数修改替代全量重训练，为大模型知识更新提供了高效途径\cite{zhang2024comprehensive}。然而，其写入机制与预训练知识的形成过程存在本质差异：预训练知识经由海量梯度迭代与多方向扰动的反复筛选，形成鲁棒的参数分布；而以ROME\cite{meng2022rome}、MEMIT\cite{meng2023memit}为代表的参数级编辑通常是一次性的局部优化，缺乏对后续参数更新的内在抵抗力。随着大模型在部署后不可避免地经历领域适配、指令对齐等持续微调，编辑知识的持久性问题逐渐显现，成为制约该范式可靠应用的关键瓶颈。

围绕编辑知识的微调衰减现象，已有研究给出了系统性证据。Cheng等人\cite{cheng2025finetune}在254个实验配置下系统量化了微调对编辑知识的影响，发现相当比例的成功编辑在后续微调后失效，且仅微调被编辑的参数层即可有效移除编辑；Wen等人\cite{wen2025retention}进一步发现，编辑知识在微调中的遗忘程度显著高于预训练固有知识，并提出改写增强与层冻结两类保留对策。上述工作均聚焦于事实性知识编辑场景，但其揭示的衰减规律引出了一个更为严峻的问题：当被编辑的对象由普通事实转变为有害知识时，编辑失效将直接导致有害能力回潮与安全护栏崩坏，安全编辑的持久性远比事实编辑场景更为紧迫。

针对知识遗忘与安全擦除的微调鲁棒性，不同技术范式给出了各自的方案。在遗忘式路线中，Wang等人\cite{wang2025ilu}将不变风险最小化引入遗忘目标，使被遗忘的知识在任意下游微调后不再回潮；Shi等人\cite{shi2025cku}通过约束知识遗忘机制，在擦除有害知识的同时保护通用能力。然而，此类方法均建立在训练式遗忘之上——遗忘本身是一个梯度下降过程，存在可施加正则约束的损失函数；而闭式参数编辑由一次矩阵运算直接写入参数，不存在可供约束的训练过程，遗忘侧的抗微调机制因而无法直接迁移至编辑范式。在安全对齐层面，Tamirisa等人\cite{tamirisa2024tar}提出的防篡改护栏使安全机制在数百步恶意微调后依然有效，但此类方法依赖重训练或大规模干预，属于模型发布前的防护手段，与知识编辑所面向的发布后轻量干预场景构成互补而非替代关系。

与此同时，攻击视角的研究从反面揭示了编辑的安全风险。Youssef等人\cite{youssef2025position}指出，知识编辑可被恶意利用以注入偏见与后门；Chen等人\cite{chen2026injection}实证了通过编辑向模型注入有害知识的可行性；Basani等人\cite{basani2026illusion}则发现，现有编辑擦除仅是脆弱的局部抑制机制，被擦除的知识在模型表征空间中结构完整，可经对抗诱导重新浮现。这些工作共同表明，编辑擦除的可靠性与持久性是真实且尚未解决的问题。

综上所述，围绕编辑与遗忘的鲁棒性已形成多条活跃研究路线：遗忘侧通过正则化约束训练过程，对齐侧依托重训练构建防篡改护栏，攻击侧持续揭示编辑擦除的脆弱性。然而，参数编辑范式下安全编辑的微调持久性研究仍处于空白。其一，现有衰减实证均基于事实编辑，安全编辑失效所引发的有害能力回潮、护栏崩坏等安全语义尚未被系统刻画；其二，更具对抗性的恶意微调场景，即攻击者以少量有害样本定向恢复被擦除知识，尚未被纳入威胁建模；其三，闭式编辑缺乏损失函数可供约束，亟需从编辑落点选择出发的新型持久化机制。与此同时，面向内容安全的编辑防御在粒度维度同样存在局限，现有方法或仅覆盖毒性单一类别\cite{lu2025toxedit}，或采用无差别的一刀切编辑强度\cite{shi2025cku}。上述空白共同指向两个核心问题：如何使参数级安全编辑在后续微调中保持有效，以及如何实现按内容风险等级精细可控的安全编辑，这正是本文两个创新点的立足点。

\subsection{现有方法存在的问题}

综合上述研究现状，面向大模型内容安全的知识编辑方法虽已取得初步进展，但仍存在以下四个亟待解决的问题。

（1）安全编辑在微调后失效且缺乏防御机制。已有实证表明，编辑知识在后续微调中会发生显著衰减\cite{cheng2025finetune,wen2025retention}，且编辑擦除仅是脆弱的局部抑制，被擦除知识可经对抗诱导重新浮现\cite{basani2026illusion}。对安全编辑而言，失效意味着被阻断的有害能力回潮与安全护栏崩坏，后果远比事实编辑失效严重；而更具对抗性的恶意微调场景尚未被纳入威胁建模。现有保留对策粒度粗且需微调方配合，参数编辑范式下尚无面向安全场景的持久化防御机制。

（2）一刀切编辑强度导致安全效用失衡。内容安全具有天然的层级结构——暴力威胁与冒犯性言论的风险等级与所需处置强度理应不同，但现有安全遗忘与编辑方法\cite{shi2025cku}均采用统一的编辑强度：强度过高则连带损伤无害知识，诱发对合法查询的过度拒绝；强度不足则高危内容漏网。安全与效用的平衡目前只能依靠全局折中，缺乏按内容风险差异化调节的能力。

（3）多类别有害知识的编辑防御覆盖不足。SafeEdit基准\cite{wang2024safeedit}已定义九类有害行为，但现有编辑防御工作多聚焦于毒性言论单一类别\cite{lu2025toxedit}，对暴力、违法建议、越狱诱导等高风险类别缺乏系统研究；不同危害类别在模型内部的知识组织与分布差异尚不明晰，致使多类别场景下的编辑定位与副作用控制缺乏依据。

（4）评估体系未覆盖持久性与分级粒度维度。现有安全编辑评估以编辑即效性（防御成功率、防御泛化性\cite{wang2024safeedit}）与过度拒绝\cite{lu2025toxedit}为主，缺少对微调后编辑保持率、按风险等级的编辑精度以及级间影响等关键维度的度量，难以全面刻画安全编辑方法在真实部署环境下的可靠性与副作用。

上述问题中，问题（1）指向安全编辑的可靠性——本文第2章围绕其研究可持久性安全编辑机制；问题（2）（3）指向安全编辑的可控性——本文第3章围绕其研究分级可控编辑框架；问题（4）贯穿两者——两章分别提出面向持久性与分级粒度的专用评估指标。

\section{本文主要内容}

本文以参数级知识编辑为技术主线，围绕内容安全场景下编辑的可靠性与可控性两大挑战，对1.2节归纳的四个问题展开研究，题目中的三个关键词在研究内容中得到逐一落实：知识编辑贯穿各项研究内容的方法基础，内容安全界定其应用场景，可控则体现为编辑效果的持久可控与编辑强度的分级可控。具体包含以下三项研究内容，总体研究框架见图~\ref{fig:chap1-framework}。

% TODO 图1-2：绘制后删除本占位框，替换为 \includegraphics[width=0.95\textwidth]{figures/chap01-fig2-framework.pdf}
\begin{figure}[htbp]
  \centering
  \fbox{\parbox[c][8.5cm][c]{0.88\textwidth}{\centering\large 图1-2 占位\\[0.6em] 本文总体研究框架\\[0.4em] \small 问题层（可靠性 / 可控性）$\rightarrow$ 方法层（第2章机制 / 第3章框架）\\ $\rightarrow$ 汇合（分级锚定，CKEdit）$\rightarrow$ 系统层（第4章防御系统）}}
  \caption{本文总体研究框架}
  \label{fig:chap1-framework}
\end{figure}

（1）面向微调鲁棒性的可持久性安全编辑机制。针对安全编辑在微调后失效且缺乏防御机制的问题，本文在参数编辑范式下系统研究安全编辑的微调衰减现象与机理。首先建立覆盖良性微调与恶意微调的双威胁模型，以衰减曲线族实证刻画安全编辑失效所引发的有害能力回潮与安全护栏崩坏等安全语义；在此基础上提出梯度投影衰减机理，并通过正交补投影干预实验完成因果验证；进而设计基于梯度正交约束的预加固编辑方法，以三重损失在编辑阶段建立对后续微调梯度的抵抗力，形成软锚定、硬锚定与预加固的方法谱系。评估方面，提出攻击回潮率、护栏完整度与新知识学习率三项指标，在SafeEdit\cite{wang2024safeedit}等基准上验证方法在双威胁场景下的持久性优势及其对通用能力的无副作用。该内容对应本文第2章。

（2）面向多类别有害知识的分级可控编辑框架。针对一刀切编辑强度导致的安全效用失衡与多类别有害知识覆盖不足问题，本文利用因果追踪技术逐层探测暴力、违法建议、自残、偏见歧视、越狱诱导等多类别有害知识在模型内部的分布，构建危害类别与模型层级的分布图谱，为分级编辑提供实证依据；进而提出基于风险感知的连续强度调度方法，按内容风险等级实现低危内容方向性衰减与高危内容彻底阻断的差异化编辑，并将分级映射设计为可替换的安全策略以适配不同部署场景；同时提出分级编辑精度与级间泄漏两项评估指标，结合分层安全效用曲线系统评估分级编辑在守住安全底线的同时降低过度拒绝的效果。该内容对应本文第3章。

（3）可控知识编辑内容安全防御系统。为将上述方法转化为可用的安全治理工具，本文设计并开发了CKEdit系统，集成知识编辑管理、分级策略配置、编辑持久监控与安全评估四大功能模块，对前两项研究内容进行系统化封装，并通过端到端集成测试验证其工程可用性。该内容对应本文第4章。

\section{本文章节安排}

全文共分五章，按照"问题——方法——系统"的主线组织，各章的具体安排如下：
\begin{enumerate}[label={（\arabic*）},leftmargin=*]
  \item 第一章为绪论。介绍本文的研究背景与意义，综述大模型知识编辑技术、生成内容安全以及知识编辑持久性与安全风险三个方面的国内外研究现状，归纳现有方法存在的四个问题，并在此基础上给出本文的主要研究内容与全文的章节安排；
  \item 第二章面向安全编辑的可靠性问题，提出面向微调鲁棒性的可持久性安全编辑机制。本章首先介绍参数级知识编辑范式，给出安全编辑问题的形式化定义，建立覆盖良性微调与恶意微调的双威胁模型；进而实证刻画安全编辑的微调衰减现象，验证基于梯度投影的衰减机理；在此基础上设计基于梯度正交约束的预加固编辑方法，形成软锚定、硬锚定与预加固的方法谱系；最后通过编辑保持率对比、消融实验与通用能力影响分析，验证所提机制在双威胁场景下的持久性优势；
  \item 第三章面向安全编辑的可控性问题，提出面向多类别有害知识的分级可控编辑框架。本章首先分析一刀切编辑导致的过度拒绝问题，给出分级可控编辑的形式化定义；进而探测多类别有害知识的层级分布图谱，为分级提供实证依据；在此基础上设计基于风险感知的连续强度调度与差异化编辑方法，并与锚定机制结合形成分级锚定；最后通过安全效用平衡分析、分级精度与级间泄漏评估以及消融与敏感性实验，系统验证框架的有效性；
  \item 第四章将前两章的方法进行系统化集成，设计并开发可控知识编辑内容安全防御系统。本章从应用场景与用户需求出发，给出系统的总体架构与工作流程，设计知识编辑管理、分级策略配置、持久监控与安全评估四大功能模块，并通过功能测试、集成效果验证与性能测试确认系统的工程可用性；
  \item 第五章为结论与建议。总结全文的主要工作与创新点，说明研究的局限性，并对后续可能的研究方向提出建议与展望。
\end{enumerate}
